\begin{examplebox}
\textbf{例1｜已知斜率直接写切线（入门秒解）}\par
求斜率为 $1$ 且与椭圆 $\frac{x^2}{4}+y^2=1$ 相切的直线。\par
\textbf{解}：$a^2=4,\ b^2=1,\ k=1$，所以
\[
m^2=4\cdot1^2+1=5,\qquad m=\pm\sqrt5.
\]
所求为
\[
\boxed{y=x+\sqrt5\quad\text{或}\quad y=x-\sqrt5.}
\]
\textbf{加深一步}：对 $m=\sqrt5$，切点为 $T(-4/\sqrt5,1/\sqrt5)$；另一条切线的切点是它关于原点的对称点。

\medskip
\textbf{例2｜过椭圆外定点引切线（保留原题并补全）}\par
过 $M(3,0)$ 引椭圆 $x^2+2y^2=2$ 的切线，求切线方程与切点。\par
\textbf{解}：椭圆化为 $\frac{x^2}{2}+y^2=1$，即 $a^2=2,b^2=1$。先设斜率存在的切线
\[
y=k(x-3)=kx-3k,
\]
故 $m=-3k$，代入相切条件：
\[
9k^2=2k^2+1\quad\Longrightarrow\quad
7k^2=1\quad\Longrightarrow\quad k=\pm\frac1{\sqrt7}.
\]
所以
\[
\boxed{y=\frac{x-3}{\sqrt7}\quad\text{或}\quad
       y=-\frac{x-3}{\sqrt7}.}
\]
对应切点分别是
\[
\boxed{\left(\frac23,-\frac{\sqrt7}{3}\right),\quad
       \left(\frac23,\frac{\sqrt7}{3}\right).}
\]
由于竖直切线仅为 $x=\pm\sqrt2$，均不经过 $M(3,0)$，故没有遗漏。

\medskip
\textbf{例3｜不求交点，直接判断位置关系}\par
对椭圆 $\frac{x^2}{9}+\frac{y^2}{4}=1$，判断斜率均为 $\frac12$、截距分别为 $2,\frac52,3$ 的三条直线与椭圆的位置关系。\par
\textbf{解}：临界值
\[
S=a^2k^2+b^2=9\cdot\frac14+4=\frac{25}{4}.
\]
\begin{itemize}
    \item $m=2$：$m^2=4<25/4$，\textbf{相交（两个交点）}；
    \item $m=5/2$：$m^2=25/4$，\textbf{相切}；
    \item $m=3$：$m^2=9>25/4$，\textbf{相离}。
\end{itemize}
\textbf{识别信号}：同一斜率、不同截距，本质是平行移动直线，无需分别联立。

\medskip
\textbf{例4｜过定点，两条切线斜率不同}\par
过 $P(1,2)$ 作椭圆 $\frac{x^2}{9}+\frac{y^2}{4}=1$ 的切线。\par
\textbf{解}：设 $y-2=k(x-1)$，则 $m=2-k$。由
\[
(2-k)^2=9k^2+4
\]
得到 $-4k=8k^2$，即 $k=0$ 或 $k=-\frac12$。故
\[
\boxed{y=2\quad\text{或}\quad y=-\frac12x+\frac52.}
\]
竖直切线为 $x=\pm3$，不经过 $P$，所以答案完整。

\medskip
\textbf{例5｜通式妙用：竖直切线也能一次判断}\par
对椭圆 $\frac{x^2}{9}+\frac{y^2}{4}=1$，判断 $2x+y=2\sqrt{10}$ 是否相切，并求切点。\par
\textbf{解}：将直线写成 $2x+y-2\sqrt{10}=0$，有 $A=2,B=1,C=-2\sqrt{10}$，于是
\[
C^2=40=9\cdot2^2+4\cdot1^2,
\]
故相切。改写为 $y=-2x+2\sqrt{10}$，利用切点公式得
\[
\boxed{T\left(\frac{9\sqrt{10}}{10},\frac{\sqrt{10}}5\right).}
\]
同理，对于 $x=3$，取 $A=1,B=0,C=-3$，也有 $C^2=9=a^2A^2+b^2B^2$，它是一条\textbf{竖直切线}。
\end{examplebox}
